% Default resume-tailor template.
% Based on the r/EngineeringResumes community template
% (https://old.reddit.com/r/EngineeringResumes/wiki/index), which itself
% derives from Jake Gutierrez's popular LaTeX resume template.
%
% This file ships as-is with the plugin. resume-tailor:generate copies it
% and fills in the candidate's real data -- it does not invent new
% preamble/macros. See skills/generate/reference.md for the structure and
% page-fit rules.

\documentclass[11pt]{article}       % set main text size
\usepackage[letterpaper,                % set paper size to letterpaper. change to a4paper for resumes outside of North America
top=0.5in,                          % specify top page margin
bottom=0.5in,                       % specify bottom page margin
left=0.5in,                         % specify left page margin
right=0.5in]{geometry}              % specify right page margin

\usepackage{XCharter}               % set font. comment this line out (or swap for \usepackage{charter}) if XCharter isn't available
\usepackage[T1]{fontenc}            % output encoding
\usepackage[utf8]{inputenc}         % input encoding
\usepackage{enumitem}               % enable lists for bullet points: itemize and \item
\usepackage[hidelinks]{hyperref}    % format hyperlinks
\usepackage{titlesec}               % enable section title customization
\raggedright                        % disable text justification
\pagestyle{empty}                   % disable page numbering

% ensure PDF output will be all-Unicode and machine-readable (ATS-friendly)
\input{glyphtounicode}
\pdfgentounicode=1

% format section headings: bolding, size, white space above and below
\titleformat{\section}{\bfseries\large}{}{0pt}{}[\vspace{1pt}\titlerule\vspace{-6.5pt}]

% format bullet points: size, white space above and below, white space between bullets
\renewcommand\labelitemi{$\vcenter{\hbox{\small$\bullet$}}$}
\setlist[itemize]{itemsep=-2pt, leftmargin=12pt, topsep=7pt} % tune these values if bullets need more/less breathing room

% resume starts here
\begin{document}

% name
\centerline{\Huge Candidate Name}

\vspace{5pt}

% contact information
\centerline{\href{mailto:name@example.com}{name@example.com} | +1 (555) 555-5555 | \href{https://github.com/username}{github.com/username}}

\vspace{-10pt}

% Optional Summary section -- most resumes don't need one (see reference.md).
% Only add when it clearly earns its space: highlighting seniority + a
% JD-matching stack for an experienced candidate, or explaining a career
% pivot. 2--3 sentences max, placed here (right after the header, before
% Skills).
%
% \section*{Summary}
% 2--3 sentence summary goes here.
%
% \vspace{-6.5pt}

% skills section
\section*{Skills}
\textbf{Languages:} Language1, Language2, Language3 \\
\textbf{Frameworks:} Framework1, Framework2 \\
\textbf{Tools:} Tool1, Tool2, Tool3

\vspace{-6.5pt}

% experience section
\section*{Experience}
\textbf{Job Title,} {Company} -- City \hfill Month Year -- Present \\
\textit{Project1 (Tech1, Tech2) --- Project2 (Tech3)} \\
\vspace{-9pt}
\begin{itemize}
  \item \textbf{XYZ (use when a real number is confirmed):} Accomplished X as measured by Y by doing Z
  \item \textbf{STAR (use for a qualitative result):} Situation Task Action Result -- past-tense action verb, no first-person pronouns
  \item \textbf{CAR (another qualitative fallback):} Challenge Action Result
\end{itemize}

\textbf{Job Title,} {Company} -- City \hfill Month Year -- Month Year \\
\textit{Project1 (Tech1, Tech2)} \\
\vspace{-9pt}
\begin{itemize}
  \item Each bullet: 1--2 lines, max one sentence. Don't let a bullet wrap onto the next line with only 1--4 words -- trim it instead
  \item Start every bullet with a strong, past-tense action verb (led, built, designed, architected, optimized...)
  \item Never fabricate a metric or a technology that isn't in the master data
\end{itemize}

\vspace{-18.5pt}

% projects section (optional -- the more work experience the candidate has, the less relevant this section becomes)
\section*{Projects}
\textbf{Project Title} \hfill \href{https://github.com/username/project}{github.com/username/project} \\
\vspace{-9pt}
\begin{itemize}
  \item Only list real projects with actual users/usage -- not mandatory school projects or tutorial-following exercises
  \item Something that solves a real problem, is actively used, and isn't abandoned in a repo
\end{itemize}

\vspace{-18.5pt}

% education section
\section*{Education}
\textbf{School} -- Degree \hfill Month Year

\end{document}
